\documentclass[11pt]{article}

\usepackage[a4paper,margin=1in]{geometry}
\usepackage{amsmath,amssymb,amsthm,mathtools}
\usepackage{booktabs}
\usepackage[hidelinks]{hyperref}

\newtheorem{theorem}{Theorem}[section]
\newtheorem{proposition}[theorem]{Proposition}
\newtheorem{lemma}[theorem]{Lemma}
\newtheorem{corollary}[theorem]{Corollary}
\theoremstyle{definition}
\newtheorem{definition}[theorem]{Definition}
\newtheorem{remark}[theorem]{Remark}

\newcommand{\R}{\mathbb{R}}
\newcommand{\Om}{\Omega}
\newcommand{\Fr}{\mathcal A}
\newcommand{\X}{\mathcal X}
\newcommand{\rstar}{\rho_*}
\newcommand{\rdel}{\rho_\delta}
\newcommand{\rhad}{\widehat\rho}
\newcommand{\Bstar}{B^*}

\title{Plan 2 / Route $\delta$: Concrete Proof Program\\
for Sharp Quantitative Faber--Krahn Stability}
\author{}
\date{}

\begin{document}
\maketitle

\begin{abstract}
This note records the current concrete route from the Plan 2
level-set machinery to the sharp quantitative Faber--Krahn inequality.
It is meant to be a compact map of the proof, not a replacement for the
individual building blocks.  The load-bearing chain is now:
\[
\begin{gathered}
\textsc{LevelSetIdentity}\to
\textsc{LimsupDeformation}\to
\textsc{SameCentreFJ}\to
\textsc{TrimmedVelocityRepair}\\
\to\textsc{WeightedMetricTrace}\to
\textsc{GoodEndpointTrace}\to
\textsc{BoundaryLayerTransfer}\to
\textsc{GlobalAssembly}.
\end{gathered}
\]
Three previously conditional links are now closed.  The conditional
\textsc{MetricFramework} hypothesis (upper-gradient recovery) is
discharged by an unconditional limsup deformation estimate
(\textsc{LimsupDeformation}); the conditional same-centre Fusco--Julin
hypothesis is discharged by a joint-functional theorem
(\textsc{SameCentreFJ}); and the exact endpoint $\rho_\delta$ at which
the metric trace is anchored can be replaced by any Borel good level
within $O(\delta_T)$ of $\rho_\delta$ (\textsc{GoodEndpointTrace}),
which preserves the sharp exponent and decouples the trace from the
exact endpoint.  The centroid and space-time notes remain useful
diagnostics, but they are not part of the shortest current proof.  The
former pointwise lower gradient hypothesis Hyp-G is also no longer a
standing assumption: on Fusco--Julin good levels, the low-gradient
part of the level surface is self-truncated and paid for by the
Cauchy--Schwarz defect $D_H$.  The remaining \emph{honest} dependencies
are the two bad-level kinetic action bounds in
\cite[Hyp.~3.6]{WeightedMetricTrace}; these were stated in
\cite{WeightedMetricTrace} as conditional inputs and remain so.
\end{abstract}

\section{Target and normalisation}

The final theorem is the sharp quantitative Faber--Krahn inequality:
there is $c_{\rm FK}(n)>0$ such that every open
$\Om\subset\R^n$, $0<|\Om|<\infty$, satisfies
\begin{equation}\label{eq:FK}
  \left(\frac{|\Om|}{|B_1|}\right)^{2/n}
  \bigl(\lambda_1(\Om)-\lambda_1(\Bstar)\bigr)
  \ge c_{\rm FK}(n)\,\Fr(\Om)^2 ,
\end{equation}
where $|\Bstar|=|\Om|$ and
\[
  \Fr(\Om):=\inf_{x_0\in\R^n}
  \frac{|\Om\Delta(x_0+\Bstar)|}{|\Bstar|}.
\]
The exponent $2$ is sharp, as shown by smooth volume-preserving
ellipsoidal perturbations of the ball.

The Plan 2 core proves first the bounded Saint--Venant stability
\begin{equation}\label{eq:SVbounded}
  \Fr(\Om)^2\le C(n,R)\,\delta_T(\Om),
  \qquad
  \delta_T(\Om):=E(\Om)-E(B_1),
\end{equation}
for $|\Om|=\omega_n$ and $\Om\subset B_R$, where
$E(\Om)=-T(\Om)/2$ and $T(\Om)$ is torsional rigidity.  The global
Faber--Krahn statement then follows from the BDV bounded reduction and
the Kohler--Jobin transfer, exactly as in the Plan 1 final library.

Throughout the bounded part we fix $|\Om|=\omega_n$ and write
$u=u_\Om$ for the torsion function, $-\Delta u=1$ in $\Om$,
$u\in H^1_0(\Om)$.  For almost every level,
\[
  E_\rho:=\{u>t(\rho)\},\qquad |E_\rho|=\omega_n\rho^n,\qquad
  d\mu(\rho):=(-t'(\rho))\,d\rho .
\]
The endpoint used in the final transfer is
\begin{equation}\label{eq:rho-delta}
  \rdel:=1-\kappa\sqrt{\delta_T(\Om)} .
\end{equation}
For small deficit, $\rdel$ lies in a fixed outer collar
$[\rho_c,1]$, for instance with $\rho_c=3/4$.  This collar localisation
is enough for the argument.

\section{Step 1: exact level-set deficit identity}

The first input is the finite-perimeter identity in
\textsc{LevelSetIdentity}.  For $E_t=\{u>t\}$, $m(t)=|E_t|$, define
\[
  \alpha(t):=\int_{\partial^*E_t}\frac{1}{|\nabla u|}\,d\mathcal H^{n-1},
  \qquad
  \beta(t):=\int_{\partial^*E_t}|\nabla u|\,d\mathcal H^{n-1},
\]
\[
  D_H(t):=\alpha(t)\beta(t)-P(E_t)^2,\qquad
  D_I(t):=P(E_t)^2-n^2\omega_n^{2/n}m(t)^{2-2/n}.
\]
Then for every open set of finite measure,
\begin{equation}\label{eq:deficit-id}
  \frac{2}{|\Om|}\bigl(E(\Om)-E(B)\bigr)
  =
  \frac{1}{|\Om|}\int_0^{\|u\|_\infty}
  \frac{D_H(t)+D_I(t)}
  {n^2\omega_n^{2/n}m(t)^{1-2/n}}\,dt .
\end{equation}
Both defects are nonnegative.  The proof uses only BV coarea on
reduced boundaries, Lipschitz truncation for the flux identity
$\int_{\partial^*E_t}|\nabla u|=|E_t|$, and the corrected rearranged
second derivative $I''=-1/\alpha$.

On every fixed collar $\rho\in[\rho_c,\rdel]$, \eqref{eq:deficit-id}
gives
\begin{equation}\label{eq:DH-DI}
  \int_{\rho_c}^{\rdel}
  \bigl(D_H(t(\rho))+D_I(t(\rho))\bigr)\,d\mu(\rho)
  \le C(n,\rho_c)\,\delta_T(\Om).
\end{equation}
We also use the Talenti comparison in two elementary forms:
$-t'(\rho)\le 1/n$, and a bad-set estimate
\[
  \bigl|\{\rho\in[\rho_c,\rdel]:-t'(\rho)<\rho_c/(2n)\}\bigr|
  \le C(n,\rho_c)\,\delta_T(\Om).
\]
The second estimate is the profile-gap input needed to pass from
$d\mu$-weighted kinetic energy to Lebesgue kinetic energy.

\section{Step 2: metric first variation}

Let $\X$ be the space of finite-measure $L^1$ classes modulo
translations:
\[
  d_\X([A],[C])=\inf_{a\in\R^n}|A\Delta(C+a)|.
\]
Set $F_\rho:=\rho^{-1}E_\rho$ so $|F_\rho|=\omega_n$.
For a Borel centre field $z_\rho$, define
\[
  V_\rho(x):=\frac{-t'(\rho)}{|\nabla u(x)|},
  \qquad
  H_{z_\rho,\rho}(x):=\frac{(x-z_\rho)\cdot\nu_\rho(x)}{\rho}
  \quad (x\in\partial^*E_\rho).
\]
The new building block \textsc{LimsupDeformation} (Theorem~1.1, Plan 2
Deliverable~A) proves the finite-perimeter metric first-variation
estimate unconditionally:
\begin{equation}\label{eq:metric-first}
  |\dot F_\rho|_\X
  \le \rho^{-n}\inf_{a\in\R^n}
  \int_{\partial^*E_\rho}
  |V_\rho-H_{z_\rho,\rho}-a\cdot\nu_\rho|\,d\mathcal H^{n-1}
\end{equation}
for almost every $\rho\in[\rho_c,1]$.  The proof is a direct coarea
computation on the symmetric difference $E_{\rho+h}\Delta T_h(E_\rho)$
with $T_h$ the affine pull-back combining homothety and a chosen
translation; no Reshetnyak or $BV$ lower semicontinuity is used.  The
previous \textsc{MetricFramework} formulation, which left
\eqref{eq:metric-first} conditional on an upper-gradient recovery
hypothesis, is therefore discharged.

\section{Step 3: isoperimetric good levels}

Fix a small dimensional threshold $\eta_I>0$ and set
\[
  G_I:=\{\rho\in[\rho_c,\rdel]:D_I(t(\rho))\le\eta_I\},
  \qquad B_I:=[\rho_c,\rdel]\setminus G_I .
\]
By \eqref{eq:DH-DI},
\begin{equation}\label{eq:BI}
  \mu(B_I)\le C\,\delta_T(\Om).
\end{equation}
On $G_I$, the building block \textsc{SameCentreFJ} (Plan 2
Deliverables~B and~C) proves the joint same-centre Fusco--Julin
package unconditionally: there is one Borel centre field
$\rho\mapsto z_\rho\in\overline{B_{R+1}}$ such that
\begin{equation}\label{eq:FJ-distance}
  d_\X(F_\rho,B_1)^2\le C(n,R,\rho_c)\,D_I(t(\rho)).
\end{equation}
The same-centre package, together with the finite-perimeter oriented
radial divergence estimate proved in
\textsc{SameCentreFJ}/\textsc{FJCenterBridge}, gives the homothetic
trace bound
\begin{equation}\label{eq:H-trace}
  \left(
    \int_{\partial^*E_\rho}|H_{z_\rho,\rho}-\bar V_\rho|\,
    d\mathcal H^{n-1}
  \right)^2
  \le C(n,R,\rho_c)\,D_I(t(\rho)),
\end{equation}
where
\[
  \bar V_\rho:=\frac1{P(E_\rho)}
  \int_{\partial^*E_\rho}V_\rho\,d\mathcal H^{n-1}
  =\frac{P(B_\rho)}{P(E_\rho)} .
\]
At this stage the same-centre theorem and its measurable-selection
input are unconditional: \cite[Thm.~3.6 + Prop.~4.2]{SameCentreFJ}.
They are no longer hidden inside \textsc{WeightedMetricTrace} and no
longer conditional on an external citation.

\section{Step 4: self-truncated velocity repair}

This is the repair replacing the former lower-gradient hypothesis.
On a good level write $S_\rho:=\partial^*E_\rho$,
$f:=|\nabla u|$, $P:=P(E_\rho)$, and
\[
  \beta:=\int_{S_\rho} f\,d\mathcal H^{n-1}=|E_\rho|,
  \qquad \bar f:=\beta/P.
\]
Split
\[
  S_\rho=G_\rho\cup L_\rho,\qquad
  G_\rho:=\{f\ge\bar f/2\},\quad L_\rho:=\{f<\bar f/2\}.
\]
The weighted variance identity for $D_H$ gives
\begin{equation}\label{eq:low-charge}
  \int_{L_\rho}\frac1f\,d\mathcal H^{n-1}
  \le \frac{4}{\beta}\,D_H(t(\rho)).
\end{equation}
Thus the low-gradient part is small in the only measure relevant to
normal velocity.  On $G_\rho$, $f\ge\bar f/2$ gives the reciprocal
variance comparison formerly supplied by a pointwise lower bound.
The resulting estimate is
\begin{equation}\label{eq:trimmed-V}
  \left(
    \int_{\partial^*E_\rho}|V_\rho-\bar V_\rho|\,d\mathcal H^{n-1}
  \right)^2
  \le C(n,\rho_c)\,D_H(t(\rho)),
  \qquad \rho\in G_I.
\end{equation}
This is the key point: tiny components, cusps, or other low-gradient
pieces do not need uniform non-degeneracy.  They are discarded at the
level surface and charged to $D_H$ itself.  The discarded swept volume
also satisfies
\[
  \int_{G_I}\int_{L_\rho}V_\rho\,d\mathcal H^{n-1}\,d\rho
  \le C(n,\rho_c)\,\delta_T(\Om).
\]

\section{Step 5: integrated action and kinetic energy}

On $G_I$, choose the same-centre field $z_\rho$ in
\eqref{eq:metric-first}, set $a=0$, and use the triangle inequality:
\[
  |V_\rho-H_{z_\rho,\rho}|
  \le |V_\rho-\bar V_\rho|+|H_{z_\rho,\rho}-\bar V_\rho|.
\]
Together with \eqref{eq:H-trace} and \eqref{eq:trimmed-V}, this gives
\begin{equation}\label{eq:metric-derivative-good}
  |\dot F_\rho|_\X^2
  \le C(n,R,\rho_c)\bigl(D_H(t(\rho))+D_I(t(\rho))\bigr)
  \qquad(\rho\in G_I).
\end{equation}
On $B_I$, the repaired \textsc{WeightedMetricTrace} no longer invokes
the unsupported perimeter-free Lipschitz bound (A0); instead it makes
its bad-level kinetic requirements explicit as conditional hypotheses.
Combining those hypotheses with \eqref{eq:DH-DI} and
\eqref{eq:FJ-distance} yields the integrated action estimate
\begin{equation}\label{eq:action}
  \int_{\rho_c}^{\rdel}
  \left(d_\X(F_\rho,B_1)^2+|\dot F_\rho|_\X^2\right)\,d\mu(\rho)
  \le C(n,R,\rho_c)\,\delta_T(\Om).
\end{equation}
Using the Talenti bad-set estimate for
$\{-t'(\rho)<\rho_c/(2n)\}$, the same bound implies the Lebesgue
kinetic estimate
\begin{equation}\label{eq:kinetic}
  \int_{\rho_c}^{\rdel}|\dot F_\rho|_\X^2\,d\rho
  \le C(n,R,\rho_c)\,\delta_T(\Om).
\end{equation}

\section{Step 6: weighted metric trace at a good endpoint}

Let $L^*:=(\rdel-\rho_c)/4$ and
$J^*:=[\rdel-L^*,\rdel]$.  From \eqref{eq:action} and the measure
lower bound $\mu(J^*)\ge c(n,\rho_c)L^*-C\delta_T$, Markov gives a
point $\rho_0\in J^*$ with
\[
  d_\X(F_{\rho_0},B_1)^2\le C\,\delta_T(\Om).
\]
For any $\rho\in J^*$, the metric absolute-continuity inequality and
\eqref{eq:kinetic} give
\[
  d_\X(F_\rho,F_{\rho_0})^2
  \le |\rho-\rho_0|\int_{J^*}|\dot F_s|_\X^2\,ds
  \le C\,\delta_T(\Om).
\]
The triangle inequality yields the uniform window trace
\[
  \sup_{\rho\in J^*}d_\X(F_\rho,B_1)^2\le C(n,R,\rho_c)\,\delta_T(\Om).
\]
The new building block \textsc{GoodEndpointTrace} (Plan 2
Deliverable~D) then promotes this to a Borel endpoint
\[
  \rhad\in[\rdel-C_1\delta_T,\rdel]\cap G_I\cap G_\tau,
  \qquad
  d_\X(F_{\rhad},B_1)^2\le C(n,R,\rho_c)\,\delta_T(\Om).
\]
The existence of $\rhad$ uses the Lebesgue bad-set measure bounds
$\mathcal L^1(B_\tau)\le C\delta_T$ (profile-gap) and
$\mathcal L^1(B_I\cap G_\tau)\le C\delta_T$ (deficit identity over
$G_\tau$).
This is exactly where the two defects interact: $D_I$ supplies many
levels close to balls, while $D_H$ supplies the kinetic control needed
to transport that closeness uniformly through the window.  The
substitution $\rdel\rightsquigarrow\rhad$ feeds the downstream
boundary-layer transfer without changing the sharp exponent (see
Step~7).

\section{Step 7: boundary-layer transfer from $\rhad$}

By definition of $d_\X$ and the scaling $F_\rho=\rho^{-1}E_\rho$,
the good-endpoint trace implies
\begin{equation}\label{eq:Erdel-asym}
  \Fr(E_{\rhad})^2\le C(n,R,\rho_c)\,\delta_T(\Om),
\end{equation}
where the asymmetry of $E_{\rhad}$ is measured at its own volume
$\omega_n\rhad^{\,n}$.  Since $E_{\rhad}\subset\Om$ and
$\rhad\in[\rdel-C_1\delta_T,\rdel]$,
\[
  |\Om\setminus E_{\rhad}|
  =\omega_n(1-\rhad^{\,n})
  \le \omega_n n(1-\rhad)
  \le \omega_n n\bigl(\kappa\sqrt{\delta_T}+C_1\delta_T\bigr)
  \le C(n)\sqrt{\delta_T(\Om)}.
\]
The elementary superlevel transfer lemma gives
\[
  \Fr(\Om)\le \Fr(E_{\rhad})
  +\frac{2}{\omega_n}|\Om\setminus E_{\rhad}|.
\]
Squaring and using \eqref{eq:Erdel-asym} gives the bounded sharp
Saint--Venant inequality
\begin{equation}\label{eq:bounded-SV}
  \Fr(\Om)^2\le C(n,R,\rho_c)\,\delta_T(\Om)
\end{equation}
in the small-deficit regime.  The cross term
$O(\delta_T^{\,3/2})$ from $\rdel-\rhad$ is dominated by the leading
$O(\sqrt{\delta_T})$ collar, so the sharp exponent is preserved.  For
large deficit the trivial bound $\Fr(\Om)<2$ is absorbed by
increasing $C$.

\section{Step 8: bounded reduction and Kohler--Jobin}

Fix $\rho_c=3/4$ for the Plan 2 core and then apply the bounded result
at the BDV dimensional radius $R_*(n)$.  The bounded-to-global
reduction imported from Plan 1 gives
\[
  E(\Om)-E(B_1)\ge c_{\rm SV}^{\rm glob}(n)\Fr(\Om)^2
  \qquad (|\Om|=\omega_n).
\]
Scaling gives the corresponding global Saint--Venant estimate at
arbitrary volume.  The Kohler--Jobin transfer then yields
\eqref{eq:FK} with
\[
  c_{\rm FK}(n)=
  \min\left\{
    \frac{4L_n\,c_{\rm SV}^{\rm glob}(n)}{(n+2)T_n},
    \frac{L_n(2^{2/(n+2)}-1)}{4}
  \right\},
\]
where $L_n=\lambda_1(B_1)$ and $T_n=T(B_1)$.

\section{What is load-bearing, and what is not}

The load-bearing proof is the eight-step chain above.  The current role
of the surrounding WIP files is as follows.

\begin{center}
\renewcommand{\arraystretch}{1.18}
\begin{tabular}{@{}lll@{}}
\toprule
File & Role & Status in the proof \\
\midrule
\texttt{WIP\_LevelSetIdentity} & exact $D_H+D_I$ identity & load-bearing \\
\texttt{WIP\_LimsupDeformation} & metric first variation, unconditional & load-bearing \\
\texttt{WIP\_SameCentreFJ} & joint same-centre FJ + Borel selection & load-bearing \\
\texttt{WIP\_MetricFramework} & smooth identity + $\X$ setup & superseded by \textsc{LimsupDeformation} \\
\texttt{WIP\_FJCenterBridge} & oriented radial divergence trace & bridge, now unconditional \\
\texttt{WIP\_TrimmedVelocityRepair} & removes Hyp-G(lower) & load-bearing \\
\texttt{WIP\_WeightedMetricTrace} & action/kinetic/uniform trace & load-bearing modulo bad-kinetic \\
\texttt{WIP\_GoodEndpointTrace} & Borel good endpoint $\rhad$ & load-bearing \\
\texttt{WIP\_BoundaryLayerTransfer} & endpoint back to $\Om$ & load-bearing \\
\texttt{WIP\_GlobalAssembly} & BDV reduction + KJ transfer & load-bearing \\
\texttt{WIP\_CentroidBound} & centroid diagnostics & auxiliary \\
\texttt{WIP\_SpaceTimeIdentity} & space-time diagnostics & auxiliary \\
\bottomrule
\end{tabular}
\end{center}

In particular, the following are not used in the current shortest
route: BDV/Cicalese--Leonardi selection, graph entry, Schauder
interpolation, nearly-spherical expansions, pointwise radial trace
(R), pointwise $(B^2)$, and the pointwise lower-gradient Hyp-G.

\section{Audit checklist}

The remaining useful audit tasks are now sharply focused:
\begin{enumerate}
\item \emph{(Discharged.)} The upper-gradient recovery theorem needed
for \eqref{eq:metric-first} is replaced by the direct finite-perimeter
limsup deformation estimate of
\cite[Thm.~1.1]{LimsupDeformation}.  Audit: verify the Lebesgue-point
input (genericity (iii) there) and the compact-strip device.
\item \emph{(Discharged.)} The exact same-centre Fusco--Julin theorem
and its measurable-selection input are now proved in
\cite[Thms.~3.6, 4.2]{SameCentreFJ}.  Audit: verify the
``centres-close'' Theorem~3.4 against the Fuglede / Cicalese--Leonardi
nearly-spherical inputs.
\item Check the self-truncated algebra
\eqref{eq:low-charge}--\eqref{eq:trimmed-V}; it is the active
replacement for Hyp-G.
\item Verify the repaired Talenti profile-gap bad-set estimate used to
pass from \eqref{eq:action} to \eqref{eq:kinetic}.
\item Verify the remaining bad-level kinetic hypotheses now stated in
\cite[Hyp.~3.6]{WeightedMetricTrace}.  These are \emph{not} discharged
by the present round of Deliverables A--D; the
\textsc{GoodEndpointTrace} variant gives a cleaner downstream
interface, but the underlying trace-uniform theorem still uses both
bad-level action bounds.
\item Check the superlevel transfer from $E_{\rhad}$ to $\Om$ and the
constant bookkeeping in the BDV/Kohler--Jobin imports.
\end{enumerate}

\begin{thebibliography}{99}

\bibitem{LevelSetIdentity}
Plan 2 building block \textsc{LevelSetIdentity},
\texttt{WIP\_LevelSetIdentity.tex}.

\bibitem{MetricFramework}
Plan 2 building block \textsc{MetricFramework},
\texttt{WIP\_MetricFramework.tex}.

\bibitem{LimsupDeformation}
Plan 2 Deliverable~A, \textsc{LimsupDeformation},
\texttt{WIP\_LimsupDeformation.tex}.

\bibitem{FJCenterBridge}
Plan 2 bridge block \textsc{FJCenterBridge},
\texttt{WIP\_FJCenterBridge.tex}.

\bibitem{SameCentreFJ}
Plan 2 Deliverables~B and~C, \textsc{SameCentreFJ},
\texttt{WIP\_SameCentreFJ.tex}.

\bibitem{TrimmedVelocityRepair}
Plan 2 building block \textsc{TrimmedVelocityRepair},
\texttt{WIP\_TrimmedVelocityRepair.tex}.

\bibitem{WeightedMetricTrace}
Plan 2 building block \textsc{WeightedMetricTrace},
\texttt{WIP\_WeightedMetricTrace.tex}.

\bibitem{GoodEndpointTrace}
Plan 2 Deliverable~D, \textsc{GoodEndpointTrace},
\texttt{WIP\_GoodEndpointTrace.tex}.

\bibitem{BoundaryLayerTransfer}
Plan 2 building block \textsc{BoundaryLayerTransfer},
\texttt{WIP\_BoundaryLayerTransfer.tex}.

\bibitem{GlobalAssembly}
Plan 2 building block \textsc{GlobalAssembly},
\texttt{WIP\_GlobalAssembly.tex}.

\bibitem{BDV-BoundedReduction-Plan1}
Plan 1 building block \textsc{BoundedReduction},
\texttt{Final/BoundedReduction.tex}.

\bibitem{KohlerJobinTransfer-Plan1}
Plan 1 building block \textsc{KohlerJobinTransfer},
\texttt{Final/KohlerJobinTransfer.tex}.

\end{thebibliography}

\end{document}
